\begin{figure}[!ht]
  \vskip -.5cm
  \centering
  %\ctikzset{
  %  amplifiers/fill=green!35,
  %}
  \begin{circuitikz}[scale=.40, rotate=90, transform shape]
    \ctikzset{
      european,
    }

    % Dimensões '\dmh' para horizontal e 'dmv' de vertical.
    \def\dmh{1.5}
    \def\dmv{3.6}
    
    % Percentual para ajuste fino.
    \def\per{.5}

    % Espaçamento.
    \def\esp{6.95}

    \foreach \a in {0,...,7} {
      \foreach \y in {0,...,3} {
        % Índices dos rótulos dos ICs.
        \pgfmathsetmacro{\z}{int(\a*4+\y+7)}

        % Índices dos rótulos dos resistores.
        \pgfmathsetmacro{\r}{int(\a*8+39+2*\y)}
        \ifnum\r>70
          \pgfmathsetmacro{\r}{int(\r+2)}
        \fi
        \pgfmathsetmacro{\s}{int(\r+1)}

        % ICXA.
        \path coordinate (NINIC\z A) at (\esp*\a,-2*\y*.8*\dmv);
        \draw (NINIC\z A) to [short] ++(\per*\dmh,0) node [op amp, anchor=+, noinv input up] (IC\z A) {IC\z A};
        % Resistores de saída de do ampop ICXA.
        \draw (IC\z A.-) to [short] ++(0,-.3*\dmv) coordinate (C1) to [short] (C1 -| IC\z A.out) to [short, -*] (IC\z A.out) to [R, l={$R_{\r}$}, a={1R}] ++(1.5*\dmh,0) coordinate (R\r);
        % Salvar coordenada globalmente.
        \path coordinate (R\r) at (R\r);

        % ICXB.
        \path coordinate (NINIC\z B) at (\esp*\a,{-(1+2*\y)*.8*\dmv});
        \draw (NINIC\z B) to [short] ++(\per*\dmh,0) node [op amp, anchor=+, noinv input up] (IC\z B) {IC\z B};
        % Resistores de saída de do ampop ICXA.
        \draw (IC\z B.-) to [short] ++(0,-.3*\dmv) coordinate (C1) to [short] (C1 -| IC\z B.out) to [short, -*] (IC\z B.out) to [R, l={$R_{\s}$}, a={1R}] ++(1.5*\dmh,0) coordinate (R\s);
        % Salvar coordenada globalmente.
        \path coordinate (R\s) at (R\s);
      }
      % Curto que conecta resistor com resistor.
      \foreach \i in {0,...,3} {
        \pgfmathsetmacro{\j}{int(39+\a*8+2*\i)}

        % Ajustar os índides para coincidir com a numeração do artigo.
        \ifnum\j>70
          \pgfmathsetmacro{\j}{int(\j+2)}
        \fi

        \pgfmathsetmacro{\k}{int(1+\j)}
        \pgfmathsetmacro{\l}{int(1+\k)}
        \draw (R\j) to [short, -*] (R\k);
        \ifnum\i<3
          \draw (R\k) to [short, -*] (R\l);
        \fi
      }
      % Conectar em paralelo o rol de ampops.
      \foreach \i in {7,...,10} {
        \pgfmathsetmacro{\j}{int(\a*4+\i)}
        \pgfmathsetmacro{\k}{int(\j+1)}
        \draw (NINIC\j A) to [short, -*] (NINIC\j B);
        \ifnum\i<10
          \draw (NINIC\j B) to [short, -*] (NINIC\k A);
        \fi
      }
    }

    % Não consegui usar o polinômio de interpolação de Lagrange.
    \draw (NINIC10B) to [R, text width=1cm, align=center, l_={$R_{37}$ 22k}] ++(0,-\dmv) node [tground] {};
    \draw (NINIC18B) to [R, text width=1cm, align=center, l_={$R_{38}$ 22k}] ++(0,-\dmv) node [tground] {};
    \draw (NINIC26B) to [R, text width=1cm, align=center, l_={$R_{71}$ 22k}] ++(0,-\dmv) node [tground] {};
    \draw (NINIC34B) to [R, text width=1cm, align=center, l_={$R_{72}$ 22k}] ++(0,-\dmv) node [tground] {};

    \foreach \i in {0,...,6} {
      \pgfmathsetmacro{\j}{int(7+4*\i)}
      \pgfmathsetmacro{\k}{int(\j+4)}

      \draw (NINIC\j A) ++(0,.5*\dmv) coordinate (CX) to [short, *-] (CX -| NINIC\k A) to [short, -*] (NINIC\k A);

      \pgfmathsetmacro{\r}{int(46+8*\i)}

       % Ajustar os índides para coincidir com a numeração do artigo.
      \ifnum\r>70
        \pgfmathsetmacro{\r}{int(\r+2)}
      \fi

      \pgfmathsetmacro{\s}{int(\r+8)}
      \draw (R\r) to [short, *-] ++ (0,-1.1*\dmv) coordinate (CX) to [short, -*] (CX -| R\s);

      \ifnum\i=6
        \draw (R\s) to [short, -o] ++(0,-1.3*\dmv) node [below=.2cm] {$K_{12}$};
      \fi
    }

    \draw (NINIC7A) to [short, *-o] ++ (0,.7*\dmv) node [above=.2cm] {$K_5$};
  \end{circuitikz}
  \vskip -.5cm
  \caption{Est\'agio de pot\^encia.}
\end{figure}
